\documentclass[11pt, letterpaper]{article}

% --- PAQUETES BÁSICOS ---
\usepackage[utf8]{inputenc}
\usepackage[T1]{fontenc}
\usepackage[spanish, es-tabla]{babel}
\usepackage{amsmath}
\usepackage{amssymb}
\usepackage{booktabs}
\usepackage{geometry}
\geometry{top=2.5cm, bottom=2.5cm, left=2.5cm, right=2.5cm}
\usepackage{hyperref}
\usepackage{xcolor}
\usepackage{listings}

% --- CONFIGURACIÓN DE LISTINGS PARA CÓDIGO PYTHON ---
\lstdefinestyle{pythonstyle}{
    language=Python,
    backgroundcolor=\color{FIMElightgray},
    commentstyle=\color{FIMEgreen}\itshape,
    keywordstyle=\color{UANLblue}\bfseries,
    numberstyle=\tiny\color{gray},
    stringstyle=\color{purple},
    basicstyle=\ttfamily\small,
    breakatwhitespace=false,
    breaklines=true,
    captionpos=b,
    keepspaces=true,
    numbers=left,
    numbersep=5pt,
    showspaces=false,
    showstringspaces=false,
    showtabs=false,
    tabsize=4
}
\lstset{style=pythonstyle}

% --- COLORES INSTITUCIONALES (FIME - UANL) ---
\definecolor{UANLblue}{HTML}{003A70}
\definecolor{UANLgold}{HTML}{FFC72C}
\definecolor{FIMEgreen}{HTML}{006A4D}
\definecolor{FIMElightgray}{HTML}{F4F4F4}

% --- CONFIGURACIÓN DE PÁGINA ---
\usepackage{fancyhdr}
\pagestyle{fancy}
\fancyhf{}
\rhead{\color{FIMEgreen}\bfseries FIME -- UANL}
\lhead{\bfseries Doctorado en Logística y Cadena de Suministro}
\rfoot{Página \thepage}
\lfoot{\bfseries Herramientas Multicriterio para Toma de Decisiones}

\begin{document}

\begin{center}
    {\LARGE\bfseries\color{UANLblue} Ejemplo Práctico Resuelto: AHP vs. BWM} \\[0.3cm]
    {\large\bfseries\color{FIMEgreen} Sesión 2: Ponderación Subjetiva de Criterios en Cadena de Suministro} \\[0.2cm]
    \rule{\linewidth}{1.5pt}
\end{center}

\section{Introducción y Planteamiento del Problema}

En el diseño y gestión de cadenas de suministro, la toma de decisiones suele ser multiatributo. Los tomadores de decisiones deben asignar pesos o niveles de importancia a criterios que muchas veces están en conflicto.

\subsection{El Escenario Logístico}
Un Centro de Distribución (CEDIS) de e-commerce ubicado en Monterrey necesita seleccionar y contratar una empresa transportista para el servicio de carga consolidada nacional (LTL). El comité directivo ha definido 4 criterios clave:
\begin{enumerate}
    \item \textbf{Costo ($C_1$):} Tarifa por kilómetro recorrido y recargos por combustible.
    \item \textbf{Confiabilidad ($C_2$):} Desempeño de entrega a tiempo y completo (OTIF).
    \item \textbf{Visibilidad ($C_3$):} Capacidad tecnológica de rastreo de embarques en tiempo real (GPS/API).
    \item \textbf{Sustentabilidad ($C_4$):} Adopción de combustibles alternos y programas de reducción de huella de carbono.
\end{enumerate}

El objetivo de este ejercicio es calcular las ponderaciones subjetivas de estos criterios mediante dos metodologías: el \textbf{Proceso de Jerarquía Analítica (AHP)} y el \textbf{Best-Worst Method (BWM)}, contrastando sus resultados y el esfuerzo requerido.

---

\section{Parte 1: Resolución mediante AHP (Saaty)}

El decisor realiza las comparaciones pareadas de los 4 criterios utilizando la escala fundamental de 1 a 9 de Saaty, resultando en la siguiente matriz de comparaciones pareadas $\mathbf{A} = [a_{ij}]$:

\begin{equation}
    \mathbf{A} = \begin{pmatrix}
        1 & 1/2 & 2 & 3 \\
        2 & 1 & 3 & 6 \\
        1/2 & 1/3 & 1 & 2 \\
        1/3 & 1/6 & 1/2 & 1
    \end{pmatrix}
\end{equation}

El número total de comparaciones requeridas fue:
\begin{equation}
    N = \frac{n(n-1)}{2} = \frac{4(3)}{2} = 6 \text{ comparaciones}
\end{equation}

\subsection{Paso 1.1: Estimación del Vector de Prioridades (Pesos)}
Para calcular el vector de prioridades aproximado, realizamos el método de normalización por columnas:
\begin{enumerate}
    \item \textbf{Suma de columnas de la matriz $\mathbf{A}$:}
    \begin{align*}
        S_{col1} &= 1 + 2 + 0.5 + 0.3333 = 3.8333 \\
        S_{col2} &= 0.5 + 1 + 0.3333 + 0.1667 = 2.0000 \\
        S_{col3} &= 2 + 3 + 1 + 0.5 = 6.5000 \\
        S_{col4} &= 3 + 6 + 2 + 1 = 12.0000
    \end{align*}
    
    \item \textbf{Normalización de cada elemento por la suma de su columna y promedio de filas:}
    Por ejemplo, para la fila 1 (Costo):
    \begin{equation*}
        w_1 \approx \frac{1}{4} \left( \frac{1}{3.8333} + \frac{0.5}{2.00} + \frac{2}{6.50} + \frac{3}{12.00} \right) \approx \frac{0.2609 + 0.2500 + 0.3077 + 0.2500}{4} \approx 0.2672
    \end{equation*}
    Realizando el mismo cálculo para las 4 filas, obtenemos el vector de pesos $\mathbf{w}$:
    \begin{equation}
        \mathbf{w} = \begin{pmatrix} w_1 \\ w_2 \\ w_3 \\ w_4 \end{pmatrix} \approx \begin{pmatrix} 0.2672 \\ 0.4958 \\ 0.1544 \\ 0.0826 \end{pmatrix}
    \end{equation}
    La suma de los pesos es $0.2672 + 0.4958 + 0.1544 + 0.0826 = 1.0000$ (perfectamente normalizada).
\end{enumerate}

\subsubsection{Implementación en Python: Cálculo de Pesos en AHP}
El cálculo anterior se automatiza en Python normalizando por columnas y promediando por filas de la siguiente manera:
\begin{lstlisting}[language=Python]
import numpy as np

# 1. Definir la matriz de comparaciones pareadas AHP
A = np.array([
    [1.0,   1/2, 2.0, 3.0],
    [2.0,   1.0, 3.0, 6.0],
    [1/2,   1/3, 1.0, 2.0],
    [1/3,   1/6, 1/2, 1.0]
])

# 2. Obtener la suma de cada columna
sum_cols = np.sum(A, axis=0)

# 3. Normalizar la matriz por columnas
norm_A = A / sum_cols

# 4. Promediar por filas para obtener los pesos finales
w_ahp = np.mean(norm_A, axis=1)

print("Pesos calculados (AHP):", w_ahp)
# Output aproximado: [0.2672, 0.4958, 0.1544, 0.0826]
\end{lstlisting}

\subsection{Paso 1.2: Análisis de Consistencia}
Para validar los resultados, debemos comprobar si el decisor ha sido razonablemente consistente.
\begin{enumerate}
    \item \textbf{Cálculo de $\mathbf{A}\mathbf{w}$:}
    \begin{equation*}
        \mathbf{A}\mathbf{w} = \begin{pmatrix}
            1 & 0.5 & 2 & 3 \\
            2 & 1 & 3 & 6 \\
            0.5 & 0.3333 & 1 & 2 \\
            0.3333 & 0.1667 & 0.5 & 1
        \end{pmatrix} \begin{pmatrix} 0.2672 \\ 0.4958 \\ 0.1544 \\ 0.0826 \end{pmatrix} \approx \begin{pmatrix} 1.0717 \\ 1.9890 \\ 0.6185 \\ 0.3315 \end{pmatrix}
    \end{equation*}
    
    \item \textbf{Elicitación\footnote{\textbf{Elicitación:} El acto de recopilar o extraer opiniones, juicios o preferencias de expertos de manera formal y estructurada.} de la estimación local de $\lambda_{\max}$:}
    Dividimos cada componente de $\mathbf{A}\mathbf{w}$ por la componente correspondiente del vector de pesos $\mathbf{w}$ y promediamos:
    \begin{align*}
        \lambda_{1} &= 1.0717 / 0.2672 \approx 4.0109 \\
        \lambda_{2} &= 1.9890 / 0.4958 \approx 4.0117 \\
        \lambda_{3} &= 0.6185 / 0.1544 \approx 4.0058 \\
        \lambda_{4} &= 0.3315 / 0.0826 \approx 4.0133 \\
        \lambda_{\max} &\approx \frac{4.0109 + 4.0117 + 4.0058 + 4.0133}{4} \approx \mathbf{4.0104}
    \end{align*}
    
    \item \textbf{Índice de Consistencia ($CI$):}
    \begin{equation}
        CI = \frac{\lambda_{\max} - n}{n-1} = \frac{4.0104 - 4}{3} \approx 0.0035
    \end{equation}
    
    \item \textbf{Razón de Consistencia ($CR$):}
    Para $n=4$, el índice aleatorio es $RI = 0.90$.
    \begin{equation}
        CR = \frac{CI}{RI} = \frac{0.0035}{0.90} \approx 0.0039 \implies \mathbf{0.39\%}
    \end{equation}
    Dado que $CR = 0.39\% < 10\%$, se concluye que los juicios de comparación pareada de AHP son extremadamente consistentes y los pesos son válidos.
\end{enumerate}

\subsubsection{Implementación en Python: Análisis de Consistencia en AHP}
Podemos verificar matemáticamente la consistencia en Python comparando el vector de pesos resultante con la matriz original utilizando el método del autovalor principal:
\begin{lstlisting}[language=Python]
# 1. Multiplicacion de la matriz por el vector de pesos (A * w)
Aw = np.dot(A, w_ahp)

# 2. Dividir cada elemento por el peso correspondiente
lambdas = Aw / w_ahp

# 3. Calcular el autovalor maximo (lambda_max)
lambda_max = np.mean(lambdas)

# 4. Calcular el Indice (CI) y Razon de Consistencia (CR)
n = A.shape[0]
CI = (lambda_max - n) / (n - 1)
RI = 0.90  # Para n = 4
CR = CI / RI

print(f"lambda_max: {lambda_max:.4f}")
print(f"Indice de Consistencia (CI): {CI:.4f}")
print(f"Razon de Consistencia (CR): {CR*100:.2f}%")
\end{lstlisting}

---

\section{Parte 2: Resolución mediante BWM (Rezaei)}

El mismo decisor resuelve el problema mediante el Best-Worst Method (BWM). 

\subsection{Paso 2.1: Identificación de Extremos (Mejor y Peor)}
\begin{itemize}
    \item \textbf{Mejor Criterio (Best, $C_B$):} Confiabilidad ($C_2$).
    \item \textbf{Peor Criterio (Worst, $C_W$):} Sustentabilidad ($C_4$).
\end{itemize}

\subsection{Paso 2.2: Vectores de Comparación Estructurada}
El decisor realiza las comparaciones en escala 1-9 del criterio Best frente a los demás, y de los demás frente al Worst:
\noindent \textbf{Vector Best-to-Others ($A_B$):}
    \begin{table}[h]
        \centering
        \begin{tabular}{ccccc}
            \toprule
            Criterios & Costo ($C_1$) & Confiabilidad ($C_2$) & Visibilidad ($C_3$) & Sustentabilidad ($C_4$) \\
            \midrule
            $C_B$ ($C_2$) & 2 & 1 & 3 & 6 \\
            \bottomrule
        \end{tabular}
    \end{table}
\noindent \textbf{Vector Others-to-Worst ($A_W$):}
    \begin{table}[h]
        \centering
        \begin{tabular}{cc}
            \toprule
            Criterios & $C_W$ ($C_4$) \\
            \midrule
            Costo ($C_1$) & 3 \\
            Confiabilidad ($C_2$) & 6 \\
            Visibilidad ($C_3$) & 2 \\
            Sustentabilidad ($C_4$) & 1 \\
            \bottomrule
        \end{tabular}
    \end{table}

El número total de comparaciones requeridas fue:
\begin{equation}
    N_{BWM} = (2n - 3) = 2(4) - 3 = 5 \text{ comparaciones}
\end{equation}
*(Se reduce en un 16.7\% el esfuerzo comparativo para solo 4 criterios; la reducción supera el 60\% para $n \ge 8$)*.

\subsection{Paso 2.3: Formulación del Modelo Lineal de Optimización}
Para calcular los pesos óptimos de BWM lineal, formulamos el siguiente Modelo de Programación Lineal:
\begin{align}
    \min \quad &\xi^L \nonumber \\
    \text{s.t.} \quad & |w_2 - 2w_1| \le \xi^L \nonumber \\
                      & |w_2 - 3w_3| \le \xi^L \nonumber \\
                      & |w_2 - 6w_4| \le \xi^L \nonumber \\
                      & |w_1 - 3w_4| \le \xi^L \nonumber \\
                      & |w_3 - 2w_4| \le \xi^L \nonumber \\
                      & w_1 + w_2 + w_3 + w_4 = 1 \nonumber \\
                      & w_j \ge 0 \quad (\forall j = 1,2,3,4)
\end{align}

\subsection{Paso 2.4: Solución Óptima del BWM}
Al resolver el sistema lineal de restricciones bajo consistencia perfecta ($\xi^{L*} = 0$):
\begin{align*}
    w_2 &= 2w_1 \implies w_1 = 0.50 w_2 \\
    w_2 &= 3w_3 \implies w_3 = 0.3333 w_2 \\
    w_2 &= 6w_4 \implies w_4 = 0.1667 w_2
\end{align*}
Sustituyendo en la restricción de suma:
\begin{align*}
    0.50 w_2 + w_2 + 0.3333 w_2 + 0.1667 w_2 &= 1 \\
    2 w_2 &= 1 \implies \mathbf{w_2^* = 0.5000}
\end{align*}
De lo cual deducimos los pesos óptimos:
\begin{equation}
    \mathbf{w}^* = \begin{pmatrix} w_1^* \\ w_2^* \\ w_3^* \\ w_4^* \end{pmatrix} = \begin{pmatrix} 0.2500 \\ 0.5000 \\ 0.1667 \\ 0.0833 \end{pmatrix}
\end{equation}
Con una inconsistencia óptima de $\xi^{L*} = 0$ (Razón de Consistencia $CR = 0.0\%$).

\subsubsection{Implementación en Python: Resolución y Consistencia del BWM}
El modelo lineal minimax de BWM se plantea y resuelve de forma declarativa y clara en Python utilizando la biblioteca de optimización PuLP:
\begin{lstlisting}[language=Python]
import pulp

# 1. Indices (0-indexed): Confiabilidad (C2) = 1, Sustentabilidad (C4) = 3
best_idx = 1
worst_idx = 3

# Vectores Best-to-Others y Others-to-Worst
a_B = [2.0, 1.0, 3.0, 6.0]
a_W = [3.0, 6.0, 2.0, 1.0]

# 2. Crear el problema de optimizacion lineal (Minimizar)
prob = pulp.LpProblem("BWM_Lineal", pulp.LpMinimize)

# 3. Definir variables de decision (pesos w_j e inconsistencia xi_L)
w = [pulp.LpVariable(f"w_{i}", lowBound=0.0, upBound=1.0) for i in range(4)]
xi = pulp.LpVariable("xi_L", lowBound=0.0)

# 4. Funcion objetivo: Minimizar xi_L
prob += xi

# 5. Restriccion de suma: w1 + w2 + w3 + w4 = 1
prob += pulp.lpSum(w) == 1.0

# 6. Restricciones de desviacion absoluta (|w_B - a_Bj * w_j| <= xi_L y |w_j - a_jW * w_W| <= xi_L)
for j in range(4):
    # Restricciones de desviacion Best-to-Others
    prob += w[best_idx] - a_B[j] * w[j] <= xi
    prob += -(w[best_idx] - a_B[j] * w[j]) <= xi
    
    # Restricciones de desviacion Others-to-Worst
    prob += w[j] - a_W[j] * w[worst_idx] <= xi
    prob += -(w[j] - a_W[j] * w[worst_idx]) <= xi

# 7. Resolver el modelo de optimizacion
prob.solve(pulp.PULP_CBC_CMD(msg=False))

# Extraer resultados
w_bwm = [pulp.value(w_var) for w_var in w]
xi_L = pulp.value(xi)

# Calcular Razon de Consistencia (CR) de BWM
tabla_CI = {1: 0.0, 2: 0.44, 3: 1.0, 4: 1.63, 5: 2.30, 6: 3.0, 7: 3.73, 8: 4.47, 9: 5.23}
a_BW = a_B[worst_idx]
CI_bwm = tabla_CI[a_BW]
CR_bwm = xi_L / CI_bwm if CI_bwm > 0.0 else 0.0

print("Pesos optimos BWM:", w_bwm)
print(f"xi_L (desviacion): {xi_L:.4f}")
print(f"CR BWM: {CR_bwm*100:.2f}%")
\end{lstlisting}

---

\section{Comparación de Resultados y Discusión Académica}

\begin{table}[h]
    \centering
    \caption{Comparativa de Pesos Obtenidos}
    \begin{tabular}{lccc}
        \toprule
        \textbf{Criterio} & \textbf{Peso AHP} & \textbf{Peso BWM} & \textbf{Diferencia Absoluta} \\
        \midrule
        Costo ($C_1$) & 0.2672 & 0.2500 & 0.0172 \\
        Confiabilidad ($C_2$) & 0.4958 & 0.5000 & 0.0042 \\
        Visibilidad ($C_3$) & 0.1544 & 0.1667 & 0.0123 \\
        Sustentabilidad ($C_4$) & 0.0826 & 0.0833 & 0.0007 \\
        \midrule
        \textbf{Comparaciones requeridas} & \textbf{6} & \textbf{5} & \textbf{-1 (16.7\% menos)} \\
        \bottomrule
    \end{tabular}
\end{table}

\subsection{Conclusiones para Investigación Doctoral}
\begin{itemize}
    \item \textbf{Convergencia:} Ambos métodos convergen a la misma estructura de preferencias: \textbf{Confiabilidad ($C_2$) $\gg$ Costo ($C_1$) $>$ Visibilidad ($C_3$) $>$ Sustentabilidad ($C_4$)}. Las discrepancias numéricas menores a 1.7\% se deben a la aproximación geométrica del autovector\footnote{\textbf{Autovector (Eigenvector):} Vector característico que en álgebra lineal no cambia su dirección bajo una transformación lineal (representada por una matriz). En AHP, representa el peso o la prioridad estable de los criterios.} en AHP.
    \item \textbf{Eficiencia:} BWM reduce significativamente las preguntas redundantes. Para un problema de 8 criterios logísticos en comités reales:
          \begin{itemize}
              \item AHP requiere $\frac{8(7)}{2} = 28$ comparaciones pareadas.
              \item BWM requiere $2(8) - 3 = 13$ comparaciones pareadas (ahorro del 53.6\%).
          \end{itemize}
    \item \textbf{Consistencia Matemática:} En BWM la formulación de optimización linealizada evita inconsistencias severas al forzar un minimax\footnote{\textbf{Minimax:} Criterio de decisión u optimización matemática que busca minimizar la máxima desviación o peor error posible en un sistema de restricciones.} que ajusta globalmente las prioridades relativas, mientras que en AHP el decisor debe ajustar valores de forma iterativa y empírica si $CR \ge 10\%$.
\end{itemize}

\section{Caso de Reto para el Estudiante: Selección de Ubicación de un nuevo CEDIS}

Con el fin de aplicar de forma autónoma los conceptos y códigos desarrollados en esta sesión, se presenta el siguiente problema de toma de decisiones estratégicas de localización.

\subsection{El Escenario del Reto}
Una empresa multinacional de retail desea abrir un nuevo Centro de Distribución (CEDIS) en la región norte de México para abastecer su canal de comercio electrónico. El comité de planificación estratégica ha seleccionado 4 criterios de localización críticos:
\begin{enumerate}
    \item \textbf{Conectividad de Infraestructura ($C_1$):} Acceso a autopistas federales, puertos secos y cercanía a vías ferroviarias.
    \item \textbf{Costos de Operación ($C_2$):} Costos de arrendamiento de la nave industrial, impuestos estatales y tarifas de servicios públicos.
    \item \textbf{Cercanía a Clientes Metropolitanos ($C_3$):} Proximidad geográfica a las principales zonas urbanas con alta densidad de demanda.
    \item \textbf{Disponibilidad de Mano de Obra ($C_4$):} Oferta y cualificación de personal técnico y operativo en logística en la zona.
\end{enumerate}

\subsection{Juicios del Decisor}
El director global de logística ha establecido las siguientes preferencias subjetivas:
\begin{itemize}
    \item \textbf{El mejor criterio (Best):} Cercanía a Clientes Metropolitanos ($C_3$).
    \item \textbf{El peor criterio (Worst):} Costos de Operación ($C_2$).
    \item \textbf{Juicios Best-to-Others ($A_B$):} El mejor criterio ($C_3$) es 2 veces más importante que Infraestructura ($C_1$), 4 veces más importante que Costos ($C_2$), y 3 veces más importante que Mano de Obra ($C_4$).
    \item \textbf{Juicios Others-to-Worst ($A_W$):} Infraestructura ($C_1$) es 3 veces más importante que Costos ($C_2$), Cercanía a Clientes ($C_3$) es 4 veces más importante que Costos ($C_2$), y Mano de Obra ($C_4$) es 2 veces más importante que Costos ($C_2$).
\end{itemize}

A partir de estas relaciones, se ha estructurado la siguiente matriz de comparación pareada de AHP:
\begin{equation}
    \mathbf{A}_{reto} = \begin{pmatrix}
        1 & 3 & 1/2 & 2 \\
        1/3 & 1 & 1/4 & 1/2 \\
        2 & 4 & 1 & 3 \\
        1/2 & 2 & 1/3 & 1
    \end{pmatrix}
\end{equation}

\subsection{Pauta de Actividades para el Estudiante}
Apoyándote en el notebook desarrollado y las explicaciones de la clase, resuelve los siguientes puntos y reporta los resultados en tu reporte de la sesión:
\begin{enumerate}
    \item \textbf{Resolución Analítica y en Python de AHP:}
          \begin{itemize}
              \item Calcula el vector de pesos aproximado de AHP normalizando la matriz $\mathbf{A}_{reto}$ por columnas.
              \item Determina el autovalor principal estimado $\lambda_{\max}$.
              \item Calcula el Índice de Consistencia ($CI$) y la Razón de Consistencia ($CR$). ¿Se cumple el umbral de aceptación del $10\%$?
          \end{itemize}
    \item \textbf{Formulación y Solución del BWM:}
          \begin{itemize}
              \item Escribe matemáticamente el modelo linealizado de optimización para este reto.
              \item Implementa y resuelve el modelo utilizando la biblioteca \texttt{PuLP} en Python.
              \item Reporta los pesos óptimos $\mathbf{w}^*$ y la desviación óptima $\xi^{L*}$.
              \item Calcula la Razón de Consistencia del BWM.
          \end{itemize}
    \item \textbf{Análisis Comparativo y Discusión:}
          \begin{itemize}
              \item Elabora una tabla que compare los pesos obtenidos por AHP y BWM.
              \item Compara el esfuerzo cognitivo (número de preguntas) que el decisor habría realizado si tuviera 8 opciones de ubicación bajo estos 4 criterios usando ambos métodos.
          \end{itemize}
\end{enumerate}

\end{document}
